\documentclass[journal,10pt]{IEEEtran}

\usepackage{amsmath,amssymb,mathtools}
\usepackage{mathrsfs}
\usepackage{amsthm}
\usepackage{booktabs}
\usepackage{multirow}
\usepackage{tabularx}
\usepackage{array}
\usepackage{algorithmic}
\usepackage{cite}
\usepackage[hidelinks]{hyperref}
\usepackage{xcolor}

\newcommand{\ac}{\mathrm{ac}}
\newcommand{\dc}{\mathrm{dc}}
\newcommand{\ph}{\varphi}
\newcommand{\cN}{\mathcal{N}}
\newcommand{\cE}{\mathcal{E}}
\newcommand{\cF}{\mathcal{F}}
\newcommand{\cB}{\mathcal{B}}
\newcommand{\cI}{\mathcal{I}}
\newcommand{\RR}{\mathbb{R}}
\newcommand{\CC}{\mathbb{C}}
\newcommand{\diag}{\operatorname{diag}}
\newcommand{\rank}{\operatorname{rank}}
\newcommand{\nullity}{\operatorname{nullity}}
\newcommand{\Schur}{\operatorname{Schur}}
\newcommand{\realify}{\mathscr{R}}
\newcommand{\TBG}{\textcolor{red!70!black}{\textbf{TBG}}}

\newtheorem{assumption}{Assumption}
\newtheorem{definition}{Definition}
\newtheorem{lemma}{Lemma}
\newtheorem{theorem}{Theorem}
\newtheorem{proposition}{Proposition}
\newtheorem{corollary}{Corollary}
\newtheorem{remark}{Remark}

\title{Constraint-Preserving Phase-Node Graph Reduction and a Boundary-Schur Interior-Point Method for Unbalanced Hybrid AC/DC Optimal Power Flow}

\author{Anonymous Authors%
\thanks{Initial manuscript draft. The theoretical development and AC-feeder reduction certificate are complete at draft level; hybrid OPF entries marked TBG must be generated before submission.}}

\markboth{IEEE Transactions Manuscript Draft}%
{Anonymous Authors: Graph-Reduced Unbalanced Hybrid AC/DC OPF}

\begin{document}
\maketitle

\begin{abstract}
This paper develops a monolithic optimal power flow (OPF) formulation for
unbalanced hybrid ac/dc distribution systems and an exact two-level reduction
method tailored to feeders with missing phases and sparse converter interfaces.
The network is represented by a compact phase-node graph containing only
physically present ac phase nodes, dc nodes, and converter coupling hyperedges.
Passive zero-injection interior phase nodes are eliminated by a
constraint-preserving Kron map, while original line-current, voltage-unbalance,
and device constraints are retained through explicit recovery operators.  We
prove a bijection between the full and reduced feasible sets, equality of their
global optimal values, correspondence of local minimizers and regular KKT
points, and a closed-form recovery formula for eliminated equality multipliers.
We further show that model-level graph reduction commutes with linearization and
with elimination of the corresponding variables from the primal-dual KKT
system.  This result enables a boundary-Schur interior-point method in which ac
feeders and dc islands are factorized independently and only converter and
controllable-device boundary variables enter the global interface solve.  A
reproducible evaluation protocol is specified for hybridized IEEE 13-, 34-,
123-, and 8500-node feeders.  An executable AC-feeder certificate demonstrates
exact boundary equivalence and sparse linear-solve acceleration; monolithic
hybrid OPF entries remain marked TBG until that solver produces auditable
results.
\end{abstract}

\begin{IEEEkeywords}
Graph reduction, hybrid ac/dc distribution system, interior-point method,
Kron reduction, missing phases, optimal power flow, Schur complement,
three-phase unbalance.
\end{IEEEkeywords}

\section{Introduction}
\IEEEPARstart{H}{ybrid} ac/dc distribution systems combine untransposed and
partially phased ac feeders, dc subnetworks, interlinking voltage-source
converters (VSCs), dc/dc converters, storage, and voltage-dependent demand.
Their OPF model is therefore larger and more tightly coupled than either a
positive-sequence hybrid OPF or an ac-only unbalanced distribution OPF.
Nevertheless, most feeder nodes are passive transit nodes, whereas nonlinear
controls and optimization variables are concentrated at a comparatively small
set of sources, loads, converters, and monitored terminals.  This separation
suggests reducing the physical graph before optimization and exploiting the
remaining block structure inside each interior-point iteration.

Three-phase unbalanced OPF has a long history
\cite{bruno2011unbalanced,rao2019helm}.  Recent formulations cover
voltage-dependent and delta-connected loads through power cones
\cite{claeys2021powercones}, explicit four-wire models and Kron-reduction error
benchmarks \cite{claeys2022fourwire}, current-voltage formulations at utility
scale \cite{soltani2024ivopf}, and meshed multiphase linearized models with
transformers and error bounds \cite{wang2026linearized}.  These works establish
that phase detail and neutral or missing-phase semantics cannot be inferred
reliably from a balanced equivalent.

Hybrid ac/dc OPF is also mature.  The \emph{hynet} framework provides a unified
hybrid formulation and convex relaxations \cite{hotz2020hynet}; distributed
methods based on ADMM and ALADIN exploit ac/dc subsystem boundaries
\cite{meyer2019distributed,zhai2022aladin}; and a recent universal branch model
supports large ac/dc systems and converter controls \cite{shahbazi2025sadra}.
For unbalanced hybrid microgrids, optimal phase and interlinking-converter power
routing has been demonstrated in simulation and hardware
\cite{esfahani2020routing}.  Accurate unbalanced hybrid power-flow methods have
also been reported using implicit Z-bus and multidimensional holomorphic
embedding \cite{pompodakis2021generic,sun2024mdhem}.  However, a general
phase-domain hybrid OPF that simultaneously addresses missing phases, exact
network reduction, and solver-level boundary condensation remains insufficiently
developed.

Kron reduction is the Schur complement of a network Laplacian or admittance
matrix and has well-established graph-theoretic properties
\cite{dorfler2013kron}.  Opti-KRON selects reductions by balancing compression
and data-fitted voltage error \cite{chevalier2022optikron}.  In four-wire OPF,
the error introduced by neutral elimination has been quantified empirically
\cite{claeys2022fourwire}.  Separately, successive Schur complements and
condensed-space KKT systems have accelerated security-constrained and GPU-based
OPF solvers \cite{kardos2020schur,shin2024gpu}.  These two uses of a Schur
complement act at different levels: Kron reduction changes the physical model
once, whereas KKT condensation eliminates Newton directions at every
optimization iteration.  Their equivalence has not, to the authors' knowledge,
been established for a constraint-preserving, missing-phase hybrid ac/dc OPF.

This paper makes the following contributions.
\begin{enumerate}
  \item A compact monolithic phase-node ac/dc OPF represents only existing ac
  phases and includes dc voltage states, phase-resolved VSC injections, dc/dc
  dispatch, converter loss and control equations, and unbalance constraints.
  \item A safe graph-reduction rule eliminates only passive zero-injection
  interior nodes.  Voltage, current, VUF, and device constraints associated with
  eliminated nodes or branches are retained as functions of boundary states.
  \item Feasible-set bijection, optimal-value equality, local-minimizer and KKT
  correspondence, second-order-condition preservation, and eliminated-dual
  recovery are proved under explicit regularity assumptions.
  \item A commutativity theorem connects model-level Kron reduction with
  linearization and KKT elimination.  It leads to a boundary-Schur IPM whose
  global system scales primarily with converters and controllable boundary
  devices rather than all feeder phase nodes.
\end{enumerate}

\begin{table*}[t]
\caption{Capability Boundary Relative to Representative Prior Work}
\label{tab:literature_position}
\centering
\footnotesize
\begin{tabular}{lllllll}
\toprule
Work & Problem & Phase detail & AC/DC & Physical reduction & KKT Schur & OPF/dual preservation\\
\midrule
Bruno \emph{et al.} \cite{bruno2011unbalanced} & OPF & unbalanced $abc$ & no & no & no & no\\
Claeys \emph{et al.} \cite{claeys2022fourwire} & OPF & four wire & no & neutral benchmark & no & no\\
Hotz and Utschick \cite{hotz2020hynet} & OPF & sequence & yes & no & no & no\\
Sun \emph{et al.} \cite{sun2024mdhem} & PF & three phase & yes & no & no & n/a\\
Kardo\v{s} \emph{et al.} \cite{kardos2020schur} & SCOPF & sequence & no & no & yes & no\\
This work & OPF & native phase masks & yes & constraint preserving & yes & primal and dual\\
\bottomrule
\end{tabular}
\end{table*}

\section{Compact Phase-Node Hybrid OPF}
\subsection{Phase-node graph}
Let $\mathcal{P}_i\subseteq\{a,b,c\}$ denote the phases physically present at
ac bus $i$.  The compact ac node set is
\begin{equation}
  \cN^{\ac}_{\phi}=\{(i,\phi):\phi\in\mathcal{P}_i\},\qquad
  n_{\phi}=\sum_i |\mathcal{P}_i|\leq 3n_{\ac}.
  \label{eq:phase_nodes}
\end{equation}
The complete hybrid graph is
\begin{equation}
  \mathcal{G}_{\phi\dc}=
  (\cN^{\ac}_{\phi}\cup\cN^{\dc},
   \cE^{\ac}\cup\cE^{\dc}\cup\cE^{c}),
\end{equation}
where $\cE^c$ contains VSC and dc/dc coupling hyperedges.  A phase-domain
admittance $Y\in\CC^{n_\phi\times n_\phi}$ includes self and mutual coupling
only among existing phase nodes.  The dc conductance matrix is
$G\in\RR^{n_{\dc}\times n_{\dc}}$.

Let $v\in\CC^{n_\phi}$ and $u\in\RR^{n_\dc}$ denote ac phase voltages and dc
voltages.  The vector $y$ collects dispatch and device variables, including
phase-resolved generation and VSC powers.  A rectangular real state is
\begin{equation}
  x=[\Re(v)^\top,\Im(v)^\top,u^\top,y^\top]^\top.
\end{equation}

\subsection{Network and converter equations}
At an ac phase node $n=(i,\phi)$, the complex balance is
\begin{equation}
 s_n^{g}(y)-s_n^{d}(v,y)+s_n^{c}(y)
 -v_n\overline{(Yv)_n}=0.                         \label{eq:ac_balance}
\end{equation}
At dc node $k$,
\begin{equation}
 p_k^{g}(y)-p_k^{d}(u,y)+p_k^{c}(y)-u_k(Gu)_k=0. \label{eq:dc_balance}
\end{equation}
For VSC $c$, with ac phases $\mathcal{P}_c$ and dc terminal $k(c)$,
\begin{align}
 p_c^{\dc}+\sum_{\phi\in\mathcal{P}_c}p_{c\phi}^{\ac}
 +\ell_c(v,u,p_c^{\ac},q_c^{\ac})&=0,             \label{eq:vsc_energy}\\
 r_c(v,u,y;\rho_c)&=0,                              \label{eq:vsc_control}
\end{align}
where $\rho_c$ specifies the converter control role.  Inequality constraints
include phase and dc voltage limits, phase-current and branch thermal limits,
VSC apparent-power, current and modulation limits, dc/dc duty limits, and
optional voltage-unbalance constraints.  For a three-phase bus,
\begin{equation}
 \mathrm{VUF}_i=\frac{|V_{2,i}|}{|V_{1,i}|}\leq\overline{\nu}_i.
 \label{eq:vuf}
\end{equation}
The full nonconvex OPF is written compactly as
\begin{subequations}\label{eq:full_opf}
\begin{align}
  \min_x\quad & f(x),\\
  \mathrm{s.t.}\quad & c(x)=0,\\
  & h(x)\leq 0,\\
  & \underline{x}\leq x\leq\overline{x}.
\end{align}
\end{subequations}
The formulation targets a regular local KKT point; no global-optimality claim is
made for the nonconvex NLP.  In the analysis below, finite variable bounds are
appended to $h$ as ordinary inequalities so that one multiplier notation covers
all inequality families.

\section{Constraint-Preserving Graph Reduction}
\subsection{Boundary and interior sets}
Partition ac phase nodes and dc nodes into retained boundary sets $\cB$ and
candidate interior sets $\cI$.  The boundary set contains every node with an
independent injection, objective term, control variable, voltage reference,
converter terminal, discrete topology decision, or nonrecoverable measurement.
An internal monitored quantity may be eliminated only if its constraint and
reported value are explicitly composed with the recovery map below.

\begin{assumption}[Exact reducibility]\label{ass:exact}
For each eliminated interior node: (i) the node is energized and has no
independent current/power injection, control, or objective term; (ii) every
constraint on its voltage or incident branch current is retained after
substitution; (iii) the interior blocks $Y_{II}$ and $G_{II}$ are nonsingular;
and (iv) ac voltage lower bounds are strictly positive, so zero complex power at
an interior node is equivalent to zero current.
\end{assumption}

With boundary-first ordering,
\begin{equation}
Y=\begin{bmatrix}Y_{BB}&Y_{BI}\\Y_{IB}&Y_{II}\end{bmatrix},\qquad
G=\begin{bmatrix}G_{BB}&G_{BI}\\G_{IB}&G_{II}\end{bmatrix}.
\end{equation}
Zero interior current gives
\begin{align}
v_I&=E_{\ac}v_B, & E_{\ac}&=-Y_{II}^{-1}Y_{IB},\label{eq:ac_recovery}\\
u_I&=E_{\dc}u_B, & E_{\dc}&=-G_{II}^{-1}G_{IB}.\label{eq:dc_recovery}
\end{align}
The reduced network matrices are
\begin{align}
Y^r&=Y_{BB}-Y_{BI}Y_{II}^{-1}Y_{IB},\label{eq:kron_ac}\\
G^r&=G_{BB}-G_{BI}G_{II}^{-1}G_{IB}.\label{eq:kron_dc}
\end{align}
For $T_{\ac}=[I\;E_{\ac}^\top]^\top$ and
$T_{\dc}=[I\;E_{\dc}^\top]^\top$, define the real recovery operator
\begin{equation}
R=\diag\!\left(\realify(T_{\ac}),T_{\dc},I_y\right),\quad
\realify(T)=\begin{bmatrix}\Re T&-\Im T\\\Im T&\Re T\end{bmatrix}.
\label{eq:R}
\end{equation}
Thus $x=Rz$ for the reduced variable $z$.

Original constraints are not discarded.  For example, the current of an
eliminated ac branch $\ell$ is
\begin{equation}
 i_\ell=H_\ell T_{\ac}v_B,\qquad
 |H_\ell T_{\ac}v_B|\leq \overline{i}_\ell,          \label{eq:recovered_current}
\end{equation}
and eliminated-bus voltage and VUF constraints are evaluated using
$v_I=E_{\ac}v_B$.  This distinguishes constraint-preserving reduction from a
boundary-equivalent model that silently removes internal operating limits.

\subsection{Feasible-set and optimal-value preservation}
Let $\cF_F$ be the feasible set of \eqref{eq:full_opf}, including passive
interior equations, and define the reduced problem by exact substitution:
\begin{subequations}\label{eq:reduced_opf}
\begin{align}
 \min_z\quad & f_R(z):=f(Rz),\\
 \mathrm{s.t.}\quad & c_R(z):=c_o(Rz)=0,\\
 & h_R(z):=h(Rz)\leq0,
\end{align}
\end{subequations}
where $c_o$ excludes the passive equations already enforced by $R$.

\begin{theorem}[Feasible-set bijection]\label{thm:feasible}
Under Assumption~\ref{ass:exact}, $R$ is a bijection from the reduced feasible
set $\cF_R$ to the full feasible set $\cF_F$.  Equivalently,
\begin{equation}
 R\cF_R=\cF_F,\qquad P_B\cF_F=\cF_R,
\end{equation}
where $P_B R=I$ is the boundary projection.
\end{theorem}

\begin{IEEEproof}
For any $z\in\cF_R$, \eqref{eq:ac_recovery} and
\eqref{eq:dc_recovery} satisfy all eliminated passive equations.  All remaining
equalities, inequalities, and bounds are precisely their full counterparts
composed with $R$; hence $Rz\in\cF_F$.  Conversely, let $x\in\cF_F$.  The
interior passive equations and nonsingularity of $Y_{II}$ and $G_{II}$ uniquely
imply \eqref{eq:ac_recovery}--\eqref{eq:dc_recovery}.  Therefore
$x=RP_Bx$, and $P_Bx$ satisfies \eqref{eq:reduced_opf}.  Uniqueness of the
interior solution makes both maps inverse on the feasible sets.
\end{IEEEproof}

\begin{corollary}[Optimality preservation]\label{cor:optimal}
The full and reduced problems have equal global optimal values.  The map $R$
gives a one-to-one correspondence between their global minimizers and between
their local minimizers.
\end{corollary}

\begin{IEEEproof}
By Theorem~\ref{thm:feasible}, feasible points correspond bijectively and
$f_R(z)=f(Rz)$.  Global optimal values and minimizers therefore correspond.
Because $R$ is linear, injective, and has continuous inverse $P_B$ on its image,
relative neighborhoods of the feasible sets also correspond, which preserves
local minimality.
\end{IEEEproof}

\subsection{KKT and dual recovery}
Write the eliminated passive equations in real form as
\begin{equation}
 a(x)=A_Bx_B+A_Ix_I=0,\qquad A_I\ \text{nonsingular},
 \label{eq:passive_real}
\end{equation}
so $E=-A_I^{-1}A_B$ and $R=[I\;E^\top]^\top$ for the corresponding state
partition.  Let $\lambda$ and $\nu$ be multipliers of $c_o$ and active rows of
$h$, respectively, and define the partial Lagrangian gradient excluding the
passive multiplier,
\begin{equation}
 q=\nabla f+J_{c_o}^{\top}\lambda+J_h^{\top}\nu
   =\begin{bmatrix}q_B\\q_I\end{bmatrix}.
\end{equation}

\begin{theorem}[KKT correspondence and eliminated-dual recovery]
\label{thm:kkt}
Suppose Assumption~\ref{ass:exact} holds and the same inequality active set is
used.  A reduced primal-dual point $(z,\lambda,\nu)$ satisfying the reduced KKT
conditions maps to a full KKT point $(Rz,\mu,\lambda,\nu)$ with
\begin{equation}
 \boxed{\mu=-A_I^{-\top}q_I}.                       \label{eq:dual_recovery}
\end{equation}
Conversely, boundary projection of any full KKT point satisfies the reduced KKT
conditions.  Under LICQ, the recovered multiplier is unique.
\end{theorem}

\begin{IEEEproof}
Reduced stationarity is $R^\top q=q_B+E^\top q_I=0$.  The internal full
stationarity row requires $q_I+A_I^\top\mu=0$, which gives
\eqref{eq:dual_recovery}.  The boundary row then becomes
$q_B+A_B^\top\mu=q_B-A_B^\top A_I^{-\top}q_I
=q_B+E^\top q_I=0$.  Primal feasibility follows from
Theorem~\ref{thm:feasible}; dual feasibility and complementarity are unchanged
because all inequalities are retained.  Conversely, premultiplying full
stationarity by $R^\top$ cancels the passive term since
$A R=0$, yielding reduced stationarity.  LICQ gives multiplier uniqueness.
\end{IEEEproof}

\begin{corollary}[Barrier-KKT correspondence]\label{cor:barrier}
For every barrier parameter $\tau>0$, append slacks $s>0$ and multipliers
$\nu>0$ so that $h(Rz)+s=0$ and $S\nu=\tau\mathbf{1}$.  The full and reduced
perturbed KKT points correspond under the same primal recovery $x=Rz$, unchanged
$(s,\nu)$, and the eliminated multiplier formula
\eqref{eq:dual_recovery}.  Hence the correspondence holds along the entire
primal--dual central path, not only after an active set is identified.
\end{corollary}

\begin{IEEEproof}
The inequality residual, slack equation, positivity requirements, and
complementarity products are identical because every original inequality is
retained as $h(Rz)$.  Only stationarity in the eliminated state differs, and
the proof of Theorem~\ref{thm:kkt} recovers its multiplier independently of
$\tau$.  Premultiplication by $R^\top$ proves the converse.
\end{IEEEproof}

\begin{corollary}[Second-order conditions]\label{cor:sosc}
If $H_L$ is the full Lagrangian Hessian, the reduced Hessian is
$H_R=R^\top H_LR$.  LICQ and the second-order necessary or sufficient condition
hold at the full point if and only if the corresponding condition holds on the
mapped reduced critical cone.
\end{corollary}

\begin{IEEEproof}
The recovery map is linear, so the chain rule gives $H_R=R^\top H_LR$ without
second derivatives of $R$.  Theorem~\ref{thm:feasible} maps feasible tangent and
critical directions bijectively through $d_x=Rd_z$, preserving the sign of the
quadratic form.
\end{IEEEproof}

\section{Commuting Reduction and Boundary-Schur IPM}
\subsection{Reduction-linearization commutativity}
The physical reduction above is performed once and is independent of the OPF
iteration.  KKT condensation instead acts on primal-dual increments at each
iteration.  The following theorem relates the two operations.

\begin{theorem}[Model/KKT Schur commutativity]\label{thm:commute}
At either a regular KKT point with fixed active set or a perturbed barrier-KKT
point with positive slacks and multipliers, forming the reduced OPF and then
linearizing its primal--dual equations produces the same Newton system as
linearizing the full OPF and exactly eliminating the passive internal state and
passive-equation multiplier increments.  In particular, up to a symmetric
permutation and congruence,
\begin{equation}
 K_R=\Schur_{(\Delta x_I,\Delta\mu)}(K_F),          \label{eq:kkt_schur_equiv}
\end{equation}
and the full primal step is recovered by $\Delta x=R\Delta z$ plus the affine
residual correction when iterates are infeasible.
\end{theorem}

\begin{IEEEproof}
Exact substitution gives the reduced residual $F_R(z)=F_o(Rz)$.  Since $R$ is
constant, $J_R=J_oR$ and $H_R=R^\top H_LR$.  Linearized passive feasibility is
$A_B\Delta x_B+A_I\Delta x_I=-a(x)$; nonsingularity of $A_I$ uniquely eliminates
$\Delta x_I$.  Internal stationarity uniquely eliminates $\Delta\mu$ by the
linearized counterpart of \eqref{eq:dual_recovery}.  Substitution into the
remaining full KKT rows yields $J_oR$ and $R^\top H_LR$, exactly the KKT blocks
obtained by differentiating the reduced problem.  Schur complementation is the
block Gaussian elimination implementing these substitutions, proving
\eqref{eq:kkt_schur_equiv}.  For barrier equations, the slack and
complementarity rows are unchanged by Corollary~\ref{cor:barrier}, so the same
elimination applies before convergence.
\end{IEEEproof}

\begin{proposition}[Missing-phase closure]\label{prop:missing}
Phase-node reduction never creates a variable for a physically absent phase.
Kron fill-in may couple retained existing phases, but the reduced node set
remains a subset of $\cN^{\ac}_{\phi}$.
\end{proposition}

\begin{IEEEproof}
Rows and columns of $Y$ are indexed only by \eqref{eq:phase_nodes}.  A Schur
complement changes entries among retained indices but cannot introduce an index
outside the original matrix.
\end{IEEEproof}

\subsection{Boundary-Schur step}
After graph reduction, partition the reduced primal-dual variables into local
ac/dc region variables $w_r$, $r=1,\ldots,m$, and global converter/control
boundary variables $w_c$.  A stabilized IPM KKT step has arrow form
\begin{equation}
\begin{bmatrix}
K_1&&&C_1^\top\\
&K_2&&C_2^\top\\
&&\ddots&\vdots\\
C_1&C_2&\cdots&K_c
\end{bmatrix}
\begin{bmatrix}\Delta w_1\\\Delta w_2\\\vdots\\\Delta w_c\end{bmatrix}
=-
\begin{bmatrix}r_1\\r_2\\\vdots\\r_c\end{bmatrix}.
\label{eq:arrow}
\end{equation}
Independent local factorizations form
\begin{align}
 S_c&=K_c-\sum_{r=1}^{m}C_rK_r^{-1}C_r^\top,\label{eq:boundary_schur}\\
 \widehat r_c&=r_c-\sum_{r=1}^{m}C_rK_r^{-1}r_r.
\end{align}
After solving $S_c\Delta w_c=-\widehat r_c$, local directions follow from
\begin{equation}
 \Delta w_r=-K_r^{-1}(r_r+C_r^\top\Delta w_c).
\end{equation}

\begin{proposition}[Boundary-step exactness]\label{prop:boundary}
If every $K_r$ and $S_c$ is nonsingular, the boundary-Schur step is identical to
the solution of \eqref{eq:arrow}.  Its arithmetic work is
\begin{equation}
 \sum_{r=1}^{m}\mathrm{fact}(K_r)
 +\mathrm{fact}(S_c)+\text{triangular solves},       \label{eq:cost}
\end{equation}
and local factorizations may execute in parallel.
\end{proposition}

\begin{IEEEproof}
The formulas are exact block Gaussian elimination of \eqref{eq:arrow}.
Equation~\eqref{eq:cost} follows by accounting for independent local
factorizations, formation and factorization of the interface matrix, and back
substitution.  Parallelism follows because no local solve depends on another
before the interface accumulation.
\end{IEEEproof}

\begin{proposition}[Sparse elimination and recovery cost]\label{prop:complexity}
Consider a sequential scalar or bus-block Kron ordering $k=1,\ldots,p$.  Let
$d_k^-$ and $d_k^+$ be the active column and row front sizes at elimination.
Without explicitly forming an inverse, numerical elimination requires
\begin{equation}
 O\!\left(\mathrm{nnz}(Y)+\sum_{k=1}^{p}d_k^-d_k^+\right)
 \label{eq:reduction_cost}
\end{equation}
arithmetic and a recovery tape of $O(\sum_k d_k^+)$ coefficients.  State
recovery has the same linear tape complexity.  A front cap and predicted-fill
cap therefore bound setup work and may leave otherwise admissible nodes
uneliminated without changing exactness for the accepted nodes.
\end{proposition}

\begin{IEEEproof}
Eliminating pivot $k$ updates each pair in the Cartesian product of its active
column and row neighbors, giving $d_k^-d_k^+$ Schur updates.  The recovery row
$x_k=-A_{kk}^{-1}A_{kN_k}x_{N_k}$ stores $d_k^+$ coefficients.  Reverse-order
substitution visits every stored coefficient once.  Rejecting a pivot changes
the retained set but not the algebra of any accepted Schur step.
\end{IEEEproof}

\subsection{Algorithm}
\begin{figure}[t]
\small
\begin{algorithmic}[1]
\REQUIRE Rich hybrid feeder model and OPF options
\STATE Build compact phase-node graph $\mathcal{G}_{\phi\dc}$
\STATE Retain injection, control, converter, constrained, and reported nodes
\STATE Classify passive zero-injection ac/dc interior nodes
\STATE Order candidates by active front size
\FOR{each candidate pivot $k$}
  \STATE Reject $k$ if singular, ill-conditioned, or above fill/front caps
  \STATE Apply its sparse Schur update and append one recovery-tape row
\ENDFOR
\STATE Form $Y^r,G^r$ and compose all original constraints with recovery
\STATE Assemble monolithic reduced $abc$--dc--VSC OPF
\FOR{each primal-dual IPM iteration}
  \STATE Evaluate reduced residuals, Jacobian, and Hessian
  \STATE Factor local region blocks $K_r$ in parallel
  \STATE Form and solve boundary system \eqref{eq:boundary_schur}
  \STATE Recover local and eliminated directions, including infeasible correction
  \STATE Filter line search; update slacks, multipliers, and barrier
\ENDFOR
\STATE Recover full states, branch flows, device dispatch, and duals
\STATE Independently evaluate the original-network KKT residual
\RETURN Full-network auditable OPF result
\end{algorithmic}
\caption{Constraint-preserving graph-reduced boundary-Schur IPM.}
\label{alg:method}
\end{figure}

\section{Numerical Evaluation and Protocol}
\subsection{Test systems and deterministic hybridization}
The evaluation uses the IEEE 13-, 34-, 123-, and 8500-node distribution
feeders in their native phase configurations.  Hybrid cases are generated by a
deterministic preprocessing script with the following rules.

\begin{table}[b]
\caption{Planned Hybrid Benchmark Configuration}
\label{tab:cases}
\centering
\footnotesize
\setlength{\tabcolsep}{3pt}
\begin{tabularx}{\columnwidth}{lccc>{\raggedright\arraybackslash}X}
\toprule
Case & Feeder & DC share & VSCs & Phase scope\\
\midrule
H13   & IEEE 13   & 15\% & 1  & Native masks\\
H34   & IEEE 34   & 20\% & 2  & Native masks\\
H123  & IEEE 123  & 25\% & 4  & Native masks\\
H8500 & IEEE 8500 & 30\% & 12 & Native masks\\
\bottomrule
\end{tabularx}
\end{table}

\begin{enumerate}
  \item Preserve all original ac phase masks, line matrices, transformers,
  regulators, loads, and source data.
  \item Rank eligible downstream subtrees by electrical distance from the
  source, breaking ties by stable source-file bus identifier.
  \item Select disjoint subtrees until a prescribed dc-load share is reached.
  Replace their internal distribution branches by dc branches with matched
  nameplate transfer capability and place one VSC at each ac/dc cut.
  \item Retain original per-phase demand on the ac side; map selected subtree
  demand to dc loads using the documented conversion efficiency.  Converter
  ratings are set to $1.25$ times the subtree peak apparent demand.
  \item Emit the full case, reduction certificate, and full-to-reduced recovery
  maps so every experiment is reproducible.
\end{enumerate}

These percentages and minimum converter counts define the initial benchmark
profile; the generated case manifest, not prose, is authoritative after case
generation.  A sensitivity sweep varies dc share, number of VSC interfaces,
loading, and the fraction of eliminable passive nodes.

\subsection{Baselines, metrics, and acceptance thresholds}
Four matched formulations isolate the two reduction levels:
\begin{enumerate}
  \item \textbf{Full-IPM}: no graph reduction and no boundary Schur solve;
  \item \textbf{GR-IPM}: graph reduction with a monolithic KKT factorization;
  \item \textbf{BS-IPM}: full graph with boundary-Schur KKT factorization;
  \item \textbf{GRBS-IPM}: both proposed levels.
\end{enumerate}
All variants use identical equations, initialization, barrier schedule,
tolerances, linear algebra backend, thread count, and stopping rules.  Reported
metrics include objective value, full-network primal and dual residuals, KKT
stationarity, recovered voltage and branch-current error, active-set agreement,
dual recovery error, variable and constraint counts, KKT nonzeros and fill,
factorization time, total time, peak memory, and speedup.

The exact-reduction implementation passes only if, after full-network recovery,
the relative objective difference is at most $10^{-8}$, the maximum state and
branch-current difference is at most $10^{-7}$ p.u., the KKT residual is at most
$10^{-7}$, and active inequalities agree within the stated complementarity
tolerance.  These are acceptance thresholds, not numerical results.

\subsection{Implemented AC-feeder reduction certificate}
Before claiming hybrid-OPF acceleration, we isolate and test the model-level
reduction on the official OpenDSS feeder scripts.  The accompanying benchmark
imports only physically present phase nodes, marks zero-independent-injection
candidates, and applies
minimum-front sequential Schur updates.  The active front is capped at 12 and
the predicted reduced-matrix nonzero count at twice the original count.  Every
accepted pivot stores a reverse recovery row.  Boundary-current equivalence is
then checked using a deterministic voltage vector, while an independently
factorized full/reduced current-injection solve checks recovered states.

Table~\ref{tab:reduction_certificate} reports a Release build on an Apple M4 Max
(16 cores, 128~GB), macOS 26.5.2, and Apple Clang 21.0.0.  Linear-solve times are
medians of seven symbolic-analysis, factorization, iterative-refinement, and
solve repetitions in one process.  The certificate eliminates 41.5--75.9\% of
the phase nodes and keeps the maximum boundary-current error below
$5.5\times10^{-8}$.  The IEEE 8500 cross-solve state difference remains
$2.2\times10^{-6}$ under the present sparse LU path despite a
$2.4\times10^{-10}$ boundary-equivalence error; it is therefore a linear-solver
conditioning diagnostic, not accepted evidence for the planned $10^{-7}$ OPF
state threshold.

% Auto-generated by tools/phase_graph_reduction_benchmark.cpp
\begin{table*}[t]
\caption{Executable AC-Feeder Phase-Node Reduction Certificate (Not Monolithic Hybrid OPF Timing)}
\label{tab:reduction_certificate}
\centering
\footnotesize
\begin{tabular}{lrrrrrrrr}
\toprule
Feeder & $n_{\phi}$ & Eliminated & Retained & $\mathrm{nnz}_r/\mathrm{nnz}$ & Reduction (ms) & Boundary error & State error & LU speedup\\
\midrule
IEEE13 & 41 & 17 (41.5\%) & 24 & 0.824 & 0.249 & $5.41\times10^{-8}$ & $2.43\times10^{-11}$ & 1.62$\times$\\
IEEE34 & 138 & 84 (60.9\%) & 54 & 0.540 & 1.188 & $2.87\times10^{-10}$ & $1.36\times10^{-13}$ & 2.90$\times$\\
IEEE123 & 278 & 161 (57.9\%) & 117 & 0.534 & 1.642 & $6.25\times10^{-9}$ & $2.71\times10^{-12}$ & 3.37$\times$\\
IEEE8500 & 11246 & 8538 (75.9\%) & 2708 & 0.333 & 38.673 & $2.32\times10^{-10}$ & $2.15\times10^{-6}$ & 3.97$\times$\\
\bottomrule
\end{tabular}
\end{table*}


The certificate validates exact sparse elimination and reverse recovery.  The
observed nodal solves are 1.6--4.1$\times$ faster.  This experiment contains
neither a dc subtree nor VSC dispatch, an objective, or a nonlinear KKT system;
it is therefore not a monolithic three-phase hybrid OPF timing result.  The
hybrid OPF fidelity and four-way IPM ablation entries below remain TBG rather
than being inferred from this narrower experiment.

% This file is intentionally separate so the benchmark generator can replace
% it without editing the manuscript. TBG means "to be generated".

\begin{table*}[t]
\caption{Full-versus-Reduced Fidelity After Recovery (All Entries TBG)}
\label{tab:fidelity}
\centering
\footnotesize
\begin{tabular}{lrrrrrr}
\toprule
Case & Eliminated phase/DC nodes & $|\Delta f|/\max(1,|f|)$ & Max. $|\Delta V|$ & Max. $|\Delta I|$ & KKT residual & Dual recovery error\\
\midrule
H13   & \TBG & \TBG & \TBG & \TBG & \TBG & \TBG\\
H34   & \TBG & \TBG & \TBG & \TBG & \TBG & \TBG\\
H123  & \TBG & \TBG & \TBG & \TBG & \TBG & \TBG\\
H8500 & \TBG & \TBG & \TBG & \TBG & \TBG & \TBG\\
\bottomrule
\end{tabular}
\end{table*}

\begin{table*}[t]
\caption{Ablation of Graph Reduction and Boundary-Schur KKT Solution (TBG)}
\label{tab:ablation}
\centering
\footnotesize
\begin{tabular}{llrrrrrr}
\toprule
Case & Method & Variables & KKT nonzeros & Fill ratio & Factorization (ms) & Total (ms) & Speedup\\
\midrule
\multirow{4}{*}{H13}   & Full-IPM & \TBG & \TBG & \TBG & \TBG & \TBG & 1.00\\
 & GR-IPM   & \TBG & \TBG & \TBG & \TBG & \TBG & \TBG\\
 & BS-IPM   & \TBG & \TBG & \TBG & \TBG & \TBG & \TBG\\
 & GRBS-IPM & \TBG & \TBG & \TBG & \TBG & \TBG & \TBG\\
\midrule
\multirow{4}{*}{H34}   & Full-IPM & \TBG & \TBG & \TBG & \TBG & \TBG & 1.00\\
 & GR-IPM   & \TBG & \TBG & \TBG & \TBG & \TBG & \TBG\\
 & BS-IPM   & \TBG & \TBG & \TBG & \TBG & \TBG & \TBG\\
 & GRBS-IPM & \TBG & \TBG & \TBG & \TBG & \TBG & \TBG\\
\midrule
\multirow{4}{*}{H123}  & Full-IPM & \TBG & \TBG & \TBG & \TBG & \TBG & 1.00\\
 & GR-IPM   & \TBG & \TBG & \TBG & \TBG & \TBG & \TBG\\
 & BS-IPM   & \TBG & \TBG & \TBG & \TBG & \TBG & \TBG\\
 & GRBS-IPM & \TBG & \TBG & \TBG & \TBG & \TBG & \TBG\\
\midrule
\multirow{4}{*}{H8500} & Full-IPM & \TBG & \TBG & \TBG & \TBG & \TBG & 1.00\\
 & GR-IPM   & \TBG & \TBG & \TBG & \TBG & \TBG & \TBG\\
 & BS-IPM   & \TBG & \TBG & \TBG & \TBG & \TBG & \TBG\\
 & GRBS-IPM & \TBG & \TBG & \TBG & \TBG & \TBG & \TBG\\
\bottomrule
\end{tabular}
\end{table*}

\begin{table}[t]
\caption{Stress and Scaling Summary (TBG)}
\label{tab:stress}
\centering
\scriptsize
\setlength{\tabcolsep}{3.5pt}
\begin{tabular}{lrrrr}
\toprule
Case & Load & DC share & VSCs & Success\\
\midrule
H13   & \TBG & \TBG & \TBG & \TBG\\
H34   & \TBG & \TBG & \TBG & \TBG\\
H123  & \TBG & \TBG & \TBG & \TBG\\
H8500 & \TBG & \TBG & \TBG & \TBG\\
\bottomrule
\end{tabular}
\end{table}


\subsection{Independent verification}
Every reduced solution is expanded and checked by an evaluator that does not
call the graph-reduction or production residual routines.  The evaluator
reconstructs the original phase-domain Y-bus and dc G-bus, recomputes all VSC
terminal balances and original line currents, and evaluates objective, primal
feasibility, complementarity, stationarity, and \eqref{eq:dual_recovery}.
Finite differences validate selected Jacobian and Hessian columns.  The ac-only
intersection is cross-checked against PowerModelsDistribution.jl
\cite{fobes2020pmd}; hybrid constraints are verified by independent equations
because common distribution tools do not provide a one-to-one dc/VSC oracle.

\section{Discussion}
\subsection{What is exact and what is not}
The theorems do not justify eliminating arbitrary PQ or voltage-dependent loads.
Such nodes remain in $\cB$ unless their nonlinear device equations and controls
are retained explicitly in the reduced problem.  Likewise, deleting original
line limits after Kron reduction changes the OPF feasible set even if boundary
voltages remain accurate.  Approximate pendant-load folding and threshold
impedance contraction may be useful engineering heuristics, but they require a
separate perturbation analysis and are outside the exact claims of this paper.

The method also does not make a nonconvex OPF globally solvable.  It preserves
the full problem's global optimum as a mathematical object and maps regular
local KKT points exactly; the IPM still returns a locally converged point unless
combined with a globally valid relaxation or enumeration method.

\subsection{Expected operating regime}
The largest benefit is expected when feeders contain many passive internal
phase nodes and comparatively few converters or controllable devices.  If
almost every node has a load, voltage constraint, or discrete control, safe
graph reduction is small.  If converters are dense, the boundary Schur system
may dominate.  Kron fill-in can also offset dimension reduction; therefore the
reduction planner must reject candidates whose predicted fill or conditioning
exceeds configurable limits.

\section{Conclusion}
This paper has formulated a unified theoretical basis for two-level reduction
of unbalanced hybrid ac/dc OPF.  A compact missing-phase graph removes phantom
variables; exact Kron maps eliminate only passive zero-injection interior
nodes; original operating constraints remain visible through recovery maps;
and a boundary-Schur IPM exploits the resulting converter-interface structure.
The feasible-set, optimal-value, KKT, second-order, dual-recovery, and
model/KKT-commutativity results establish when the reduced calculation is
mathematically the same optimization problem rather than an approximate grid
equivalent.  Completing the reproducible H13, H34, H123, and H8500 experiments
and replacing all TBG entries are required before submission.

\IEEEtriggeratref{10}
\bibliographystyle{IEEEtran}
\bibliography{references}

\end{document}
